\documentclass[a4paper]{article}
\usepackage{amsfonts}
\usepackage{amsmath}
\usepackage{amssymb}
\usepackage{graphicx}     

\begin{document}
  
\noindent \large Auteur: Abdellah El Moussaoui \\
\noindent \large Studierichting: Informatica\\
\noindent \large Oefeningenreeks 1D nr. 19.6\\

\medskip

\normalsize

Bepaal al de nulpunten van de volgende veelterm in C en ontbind de veelterm in factoren over R[z]: \\

$z^3 + z - 2$\\

\begin{tabular}{c|cccc}
& $1$ & $0$ & $1$ & $-2$\\
$1$ & & $1$ & $1$ & $2$\\ \hline
& $1$ & $1$ & $2$ & $0$
\end{tabular}\\

$(z - 1)(z^2 + z + 2)$\\

$D = 1^2 - 4(1)(2) = -7 $\\

$z_2 = \dfrac{-1+i\sqrt{7}}{2}$ , $z_3 = \dfrac{-1-i\sqrt{7}}{2}$\\

$\Rightarrow \left(z-1\right)\left(z+\dfrac{1}{2}-\dfrac{\sqrt{7}}{2}i\right)\left(z+\dfrac{1}{2}+\dfrac{\sqrt{7}}{2}i\right)$\\\\

Dus de nulpunten zijn:\\

$\boxed{z_1 = 1\ ,\ z_2 = \dfrac{-1+i\sqrt{7}}{2}\ ,\ z_3 = \dfrac{-1-i\sqrt{7}}{2}}$\\


\begin{thebibliography}{999}
%\bibitem{a} Referentie
\end{thebibliography}
\end{document} 
